\section{Metric slices and real stability}\label{sec:slices}

We first identify the complex quotients at arbitrary real parameters.
This also gives a properness estimate for the whole family, which will
be needed for intrinsic metric convergence. We then prove a quantitative
slice theorem at an arbitrary singular point.

\subsection{Balancing at arbitrary real parameters}
For a quiver representation $x$ of dimension $m$, a collection
$D=(D_i\subset W_i)$ is a subrepresentation if
$x_aD_{t(a)}\subset D_{h(a)}$ for every arrow.
Let $\eta\in\mathbb R^{Q_0}$ satisfy $\eta\cdot m=0$.
We call $x$ $\eta$-semistable if $\eta\cdot\dim D\ge0$ for every
subrepresentation, and $\eta$-stable if the inequality is strict for
every nonzero proper subrepresentation. It is $\eta$-polystable if
it is a direct sum of $\eta$-stable representations, each of total
$\eta$-weight zero. These definitions use arbitrary real $\eta$.

\begin{proposition}\label{prop:real-polystability}
An orbit $Gx$ meets $\mu^{-1}(\eta)$ if and only if $x$ is
$\eta$-polystable. The intersection is one $K$-orbit.
For a balanced point, $G_x$ is the complexification of $K_x$.
A stable representation has only scalar endomorphisms.
\end{proposition}
\begin{proof}
For stable $x$, consider the left $K$-invariant function
\[
 \mathcal F(g)=\|gx\|^2-2\sum_i\eta_i\log|\det g_i|.
\]
Multiplying by a common positive scalar does not change it.
Polar decomposition therefore reduces minimization to
\[
 F(S)=\|e^Sx\|^2-2\langle\eta,S\rangle,\qquad
 S\in\mathfrak a,\quad \operatorname{tr}S=0.
\]
This function is proper. Indeed, let $\|S_n\|=r_n\to\infty$ and
pass to a subsequence with $S_n/r_n\to T$, $\|T\|=1$.
Diagonalize vertexwise and pass to convergent unitary eigenbases and
normalized eigenvalues. If a nonzero coefficient of $x$ in the limiting
bases goes from a $T$-eigenvalue $\tau$ to a larger $\tau'$, its
contribution to $\|e^{S_n}x\|^2$ grows at least exponentially in $r_n$,
whereas the trace term is linear.

Otherwise every cumulative lower eigenspace of $T$ is invariant.
For its distinct global eigenvalues $\tau_1<\cdots<\tau_l$, let
$D_j=W_{\le\tau_j}$ and let $P_j$ be the corresponding projection.
There are at least two distinct eigenvalues, since $T$ is nonzero
and trace-free. Thus
\[
 T=\tau_l I-\sum_{j<l}(\tau_{j+1}-\tau_j)P_j,\qquad
 \langle\eta,T\rangle
 =-\sum_{j<l}(\tau_{j+1}-\tau_j)\eta\cdot\dim D_j<0.
\]
It follows that $F(S_n)\ge-2r_n\langle\eta,S_n/r_n\rangle\to\infty$.
These subsequence alternatives prove properness. A minimum exists,
and varying $\mathcal F$ by $g\mapsto e^{tA}g$ gives
$2\langle\mu(gx)-\eta,A\rangle=0$ for every Hermitian $A$.
The scalar equation vanishes identically, so $\mu(gx)=\eta$.

For a balanced representation and an invariant $D$, compression of the
moment equations by the orthogonal projection $P_D$ gives
\begin{equation}\label{eq:subrepresentation-trace}
 \eta\cdot\dim D
 =\sum_a\|P_{D_{h(a)}}x_a(1-P_{D_{t(a)}})\|^2\ge0.
\end{equation}
Equality means that $D$ is reducing. Split orthogonally whenever
equality occurs for a proper nonzero $D$ and induct on total dimension.
This gives stable summands of zero total weight. Conversely balance
each stable summand by the preceding minimization and take their
orthogonal direct sum.

Uniqueness follows from polar decomposition and convexity as in
Theorem~\ref{thm:uniform-balancing}: for balanced $x,e^Sx$ the function
$\|e^{tS}x\|^2-2t\langle\eta,S\rangle$ has zero endpoint derivatives
and second derivative $4\|C_{e^{tS}x}S\|^2$, so $e^Sx=x$.
If $g$ stabilizes $x$, its compact and positive polar factors
therefore both stabilize $x$, proving $G_x=(K_x)^{\mathbb C}$.

Finally, a nonzero singular endomorphism of a stable representation
would have nonzero proper kernel and image, both of positive weight,
while their weights add to zero. Every nonzero endomorphism is
therefore invertible. Subtracting an eigenvalue proves it scalar.
\end{proof}

The proposition is the polystable form of the real Kempf--Ness
criterion. Its direct proof avoids any genericity restriction.
In particular, a balanced orbit with a positive-dimensional stabilizer
is not incorrectly declared stable.

\subsection{Invariant properness and fixed faces}
From now on $V$ is a doubled-quiver representation, with
$Z=\mu_{\mathbb C}^{-1}(0)$ given its reduced structure.
Choose positive-degree homogeneous generators $f_1,\ldots,f_q$ of
$\mathbb C[Z]^G$, of degrees $d_j$, and set
$a(x)=\max_j|f_j(x)|^{1/d_j}$; set $a=0$ if there are no generators.
Finite generation is the standard reductive invariant-theory theorem
\cite[\S1.6]{KP96}.

\begin{lemma}\label{lem:invariant-properness}
There is a constant $C$ such that
\begin{equation}\label{eq:invariant-properness}
 \|x\|\le C\bigl(\|\mu(x)\|^{1/2}+a(x)\bigr),\qquad x\in Z.
\end{equation}
Consequently the family of balanced compact quotients maps properly
to the product of the central real parameter space and $Z/\!/G$.
Representatives in the affine fiber over the origin satisfy
$\|x\|\le C\sqrt{\|\eta\|}$.
\end{lemma}
\begin{proof}
The right-hand parenthesis is strictly positive on the unit sphere
in $Z$. A zero there would be balanced at zero and have all invariant
coordinates equal to those of the origin. Zero-character Kempf--Ness
makes its complex orbit closed, and the affine quotient fiber has
the origin as its unique closed orbit
\cite[\S2 and Theorem~6.1]{King94}, a contradiction.
Its positive minimum and real homogeneity prove the estimate.
Bounded parameter and invariant coordinates then bound all balanced
representatives. Their inverse image of a compact set is closed
before quotienting, hence compact. The last assertion follows when
all $f_j$ vanish.
\end{proof}

Partition $\Theta=\{\eta:\eta\cdot m=0\}$ by the signs of $\eta\cdot d$
for every integral $0\le d\le m$, excluding $0,m$. Each relatively
open face $F$ is rational polyhedral and contains a rational parameter;
clear denominators to choose an integral $\theta_F\in F$.
Our character is
$\chi_{\theta_F}(g)=\prod_i\det(g_i)^{(\theta_F)_i}$.
This is King's positive subrepresentation convention
\cite[Propositions~3.1 and 3.2]{King94}.
All stable, semistable and polystable objects are constant on $F$:
the sign list includes dimensions of every possible summand and
its subrepresentations.

\begin{corollary}\label{cor:fixed-face}
For $\eta\in F$, the fixed algebraic quotient map induces a
homeomorphism
\[
 M_\eta=(Z\cap\mu^{-1}(\eta))/K
       \ \longrightarrow\ Q_F=Z/\!/_{\chi_{\theta_F}}G.
\]
\end{corollary}
\begin{proof}
King's description of the closed semistable orbits and
Proposition~\ref{prop:real-polystability} give a continuous bijection.
Over a compact subset of $Q_F$, affine invariant coordinates are
bounded, so Lemma~\ref{lem:invariant-properness} bounds representatives.
Every convergent sequence of them has an $\eta$-balanced limit,
which is still $\theta_F$-polystable. The quotient map is therefore
defined at the limit and takes it into the same compact set.
The bijection is proper, hence a homeomorphism.
\end{proof}
We shall use the fixed algebraic quotient maps on representative
spaces. The corollary does not by itself give a jointly continuous
inverse as the real parameter varies.

\subsection{A quantitative Hamiltonian slice}
Fix $x_*$ with $\mu_{\mathbb C}(x_*)=0$ and $\mu(x_*)=\eta_*$
central, and let $H=K_{x_*}$.
The quaternionic normal representation is
\[
 W=(L\oplus JL)^\perp,\qquad L=T_{x_*}(Gx_*),
\]
where $J$ is the antiunitary quaternionic structure on $V$.
The complex orbit tangent $L$ is holomorphically symplectic isotropic,
$L\perp JL$, and $L\oplus JL$ is quaternionic. Thus $W$ inherits
its Hermitian and holomorphic symplectic forms.

\begin{theorem}\label{thm:metric-slice}
There are nested common neighborhoods, for $\eta$ near $\eta_*$,
identifying $M_\eta$ near $[x_*]$ with the quotient of the full normal
representation $W$ at central value $\eta|_{\mathfrak h}$.
Here $\mathfrak h$ denotes the Hermitian stabilizer and restriction
is dual to its actual embedding in the vertex algebra.
The identifications are homeomorphisms onto open neighborhoods,
smooth on strata, and their quotient metrics satisfy
\begin{equation}\label{eq:slice-metric}
 (1-C(|\eta-\eta_*|+\|w\|))g^W_{\eta|_{\mathfrak h}}
 \le f_\eta^*g_\eta
 \le(1+C(|\eta-\eta_*|+\|w\|))g^W_{\eta|_{\mathfrak h}}.
\end{equation}
Constants depend on the fixed center $x_*$, but not on positive
eigenvalues of the retained $H$-orbit operator at $w$.
The representation $W$ is unitarily a doubled Ext-quiver representation,
including paired loops.
\end{theorem}
\begin{proof}
Shift the real moment by $\eta_*$ and apply the full holomorphic
Hamiltonian slice theorem of Mayrand
\cite[Theorem~1.3, Proposition~3.1 and (3.5)]{May22}.
The ambient model is
\[
 G\times_{H^{\mathbb C}}(\mathfrak h^\circ\times W),
 \qquad
 [g,\xi,w]\longmapsto
 \operatorname{Ad}_g^*(\xi+\mu_{\mathbb C,W}(w)).
\]
The two target summands are complementary, so the complex zero
locus is exactly $\xi=0$, $\mu_{\mathbb C,W}(w)=0$.
The theorem applies in the shifted semistable open neighborhood;
closedness of $Gx_*$ in the original affine space is unnecessary.

We need the slice differential to be the isometric normal inclusion.
Identify the orbit directions with $L$, their symplectic cotangent dual
with $JL$ by the fixed orbit Gram pairing, and the normal with the
actual $W$ by the identity. This is an $H$-equivariant symplectic
linear identification of full tangents, isometric on $W$.
It can be prescribed in the holomorphic slice: the derivative of an
initial slice differs by $(a,n)\mapsto(a+An,Bn)$, and the local
map $[g,n]\mapsto[g\exp(An),Bn]$ realizes this change
\cite[Theorem~3.9 and \S3.5]{May22}.
The pulled-back and model symplectic forms now agree as full tensors
along the zero orbit. Relative radial homotopy produces a primitive
of their difference vanishing to second order in normal variables;
the equivariant Moser vector field has the same order. Its flow has
identity first jet on the zero section and preserves the prescribed
normal isometry. This is the relative Moser construction of
\cite[Theorem~3.3 and Remark~3.4]{May22}, with the order of vanishing
made explicit. Restricting to $[1,0,w]$ gives a full-ball holomorphic
map $F$ with $F(0)=x_*$ and $DF(0)$ the isometric inclusion.

For the representation claim, decompose the finite-dimensional star
algebra generated by vertex projections, arrows of $x_*$ and their
adjoints:
\[
 E_i=\bigoplus_a D_i^a\otimes N_a,\qquad
 (x_*)_e=\bigoplus_a(x^a)_e\otimes I_{N_a},\qquad
 H=\prod_aU(N_a).
\]
Proposition~\ref{prop:real-polystability} identifies its complexification
with the full complex stabilizer. The factors are stable objects at
$\eta_*$, not necessarily simple modules.
The moment deformation complex splits into factor complexes tensored
with $\operatorname{Hom}(N_a,N_b)$. Its orthogonal first cohomology is
\[
 W=\bigoplus_{a,b}\mathcal E_{ab}\otimes\operatorname{Hom}(N_a,N_b),
\]
with product Hermitian norm. Quaternionic duality pairs
$\mathcal E_{ab}$ with $\mathcal E_{ba}$ antiunitarily and makes each
$\mathcal E_{aa}$ quaternionic. Paired orthonormal bases give the
unitary doubled Ext quiver, including pairs of loops. This is the
metric realization of the usual Ext-quiver slice
\cite[\S3.2]{BeSc24}. If $d_i^a=\dim D_i^a$, restriction of a central
covector gives precisely
\begin{equation}\label{eq:normal-level}
 \zeta_a=\sum_i d_i^a\eta_i.
\end{equation}

Let $\mathfrak q=\mathfrak h^\perp$ be the fixed Hermitian complement.
Solve
\[
 P_{\mathfrak q}\bigl(\mu(e^TF(w))-\eta\bigr)=0,\qquad
 T\in\mathfrak q.
\]
Its derivative at $(\eta_*,0,0)$ has positive quadratic form
$2\|C_{x_*}T\|^2$, invertible on $\mathfrak q$.
The implicit-function theorem gives a unique small equivariant
$\sigma(\eta,w)$ with
$\sigma=O(|\eta-\eta_*|+\|w\|^2)$; the linear $w$ term vanishes
by normal orthogonality. Put $\Psi_\eta=e^\sigma F$.
At $w=0$ this is $H$-fixed. The original homogeneous moment paired
with $\mathfrak h$ vanishes on the entire $H$-fixed subspace.
Also $\eta_*|_{\mathfrak h}=0$, since
$\langle\mu(x_*),S\rangle=\operatorname{Re}\langle C_{x_*}S,x_*\rangle$.
Thus the remaining moment $\rho_\eta=\mu(\Psi_\eta)|_{\mathfrak h}$
is normalized at zero, with target
$\eta|_{\mathfrak h}=(\eta-\eta_*)|_{\mathfrak h}$.

The form $\Omega_\eta=\Psi_\eta^*\omega_0$ is closed and invariant,
has this exact moment map, and differs from the flat form on $W$ by
$O(|\eta-\eta_*|+\|w\|)$. Apply
Theorem~\ref{thm:uniform-balancing} to the full retained $H$ quiver.
This gives both directions of the finite correction uniformly in its
weak orbit eigenvalues.

For the metric, transport the ambient Hermitian norm through
$e^\sigma$ and define on the full tangent space to $W$
\[
 B_{\eta,w}(v,v)=
 \inf_{T\in\mathfrak q^{\mathbb C}}
       \|DF(w)v-C_{F(w)}T\|_{e^\sigma}^2.
\]
The complementary Gram matrix is invertible at the center and nearby,
so $B$ is smooth, equals the flat normal norm at the center, and has
error $O(|\eta-\eta_*|+\|w\|)$.
Exact equivariance of $F$ and splitting of the full orbit algebra give
\[
 \inf_{\gamma\in\mathfrak g}
       \|DF(w)v-C_{F(w)}\gamma\|_{e^\sigma}^2
 =\inf_{\xi\in\mathfrak h^{\mathbb C}}
       B_{\eta,w}(v-C_w\xi,v-C_w\xi).
\]
On genuine integrability-stratum tangents the left side is the
quotient norm. Derivatives of $\sigma$ and of the retained complex
correction are orbit directions and disappear in this infimum.
The correction has norm $O(|\eta-\eta_*|+\|w\|)$, so its action
distorts squared norms by $1+O(|\eta-\eta_*|+\|w\|)$.
This proves \eqref{eq:slice-metric}.

For clarity, smoothness on a fixed stratum follows from a separate
implicit-function argument. At a balanced point with compact
stabilizer $J$, the correction commutes with $J$, and its complex
stabilizer remains $J^{\mathbb C}$. On the fixed subspace $W^J$ the
remaining moment takes values in the commutant of $J$ orthogonal to
its Hermitian algebra. Its derivative in that positive-gauge space has
quadratic form $2\Omega(iC_yT,C_yT)>0$ for $T\ne0$.
The implicit-function theorem gives a smooth local correction there.
A compact orbit-type slice and uniqueness identify its local maps
on overlaps. Only the inverse at the fixed point is used.

Finally the full complex Hamiltonian slice classifies every nearby
integrable orbit. Write its small complex gauge in polar form. Its
positive logarithm can be normalized to lie in $\mathfrak q$ by
varying the positive $H$ factor: if the original logarithm is $s$,
the derivative at zero of the $\mathfrak h$ projection of
\[
 \frac12\log(e^{-h}e^{2s}e^{-h})
\]
with respect to $h\in\mathfrak h$ is $-I$.
The finite implicit-function theorem supplies this normalization.
Uniqueness of the complementary solution identifies it with $\sigma$;
the reverse part of Theorem~\ref{thm:uniform-balancing} then gives a
flat normal-balanced representative. All constructions hold on common
nested parameter neighborhoods. Compact-quotient uniqueness proves
continuity of the inverse, giving full local coverage.
\end{proof}

Write $U=W^H$ and $W=U\oplus W'$ orthogonally. These are symplectic
summands; the moments ignore $U$. The normal quotient is therefore
the exact metric product
\begin{equation}\label{eq:normal-product}
 W///_\zeta H=U\times(W'///_\zeta H).
\end{equation}
The stratum through the center is $U$ times the normal vertex.
In particular, holding a balanced $z\in W'$ fixed while following a
path in $U$ has exactly its Euclidean $U$-length in the model.
The full $W$ is a doubled-quiver representation; $W'$ alone may
contain traceless loop spaces and need not be one.

